\chapter{The MDL Growth Law}\label{ch:mdl}

The question this chapter answers: \emph{when is a new concept worth
inventing?} Earlier designs answered it with a threshold someone chose.
This one derives a threshold instead.

\section{Two-part description length}

\begin{intu}\Intu
A memory is a compression of experience. Minimum description length
makes that literal: the total cost of an explanation is what it takes to
write down the model plus what it takes to write down the data
\emph{given} the model. Adding a concept always makes the first term
worse. It is worth doing exactly when it makes the second term better by
more.
\end{intu}

\begin{definition}[Two-part MDL]\label{def:mdl}
\[
  L_{\text{total}}
  =\underbrace{L(\text{model})}_{\text{the memory}}
  +\underbrace{L(\text{data}\mid\text{model})}_{\text{the experience}} .
\]
The \emph{model} is the complex $\CC$ together with the sheaf $\Fs$ ---
what the system \emph{is}. The \emph{data} is the observed record: the
sequence of assemblies and the succession relation read off it --- what
the system has \emph{seen}.
\end{definition}

Neither term alone is the criterion. Minimising the model term alone
forbids all growth; minimising the data term alone is achieved by
memorising everything. The criterion is the sum.

\section{The model term}

\begin{proposition}[Cost of coning a cycle]\label{prop:modelterm}
\Proved\ Let $\gamma$ be a $1$-cycle of length $k$, and cone it with a
new vertex $w$ carrying a stalk of dimension $d_w$. Writing $b$ for bits
per real number, the added model cost is
\begin{equation}\label{eq:modelterm}
  \Delta L(\text{model})
  =b\,d_w\Bigl(1+\sum_{e\in\gamma}d_e\Bigr).
\end{equation}
\end{proposition}
\begin{proof}
Three kinds of new object are described. The vertex with its stalk costs
$d_w\cdot b$ bits. Each of the $k$ restriction maps out of $\stalk{w}$
is a $d_e\times d_w$ matrix, costing $d_w d_e b$ bits, summed over
$e\in\gamma$. The $k$ new triangles cost \emph{nothing}: they are
determined by the cone construction, not described separately. Summing
gives \eqref{eq:modelterm}.
\end{proof}

\begin{remark}[Two consequences, both load-bearing]
First, the cost is \emph{linear in $d_w$}. The dimension of the new
stalk is the dominant free quantity, so whatever fixes $d_w$ fixes most
of the cost --- which is what Section~\ref{sec:conesize} is for.
Second, \emph{triangles are free}. They are implied by the cone, so the
$k$ new $2$-cells cost nothing to describe. This is a real asymmetry,
not a bookkeeping convenience, and it is why coning is cheap relative to
what it buys.
\end{remark}

\section{The data term and the law}

\begin{proposition}[Data-term saving]\label{prop:dataterm}
\Proved\ Let the cycle $\gamma$ occur $n$ times in the record. Without
the concept, each traversal is coded as $k$ separate transitions, each
costing its full surprisal; with the concept it is a single named event.
Then
\begin{equation}\label{eq:dataterm}
  \Delta L(\text{data}\mid\text{model})=-n\bigl(c_{\text{old}}-c_{\text{new}}\bigr),
\end{equation}
where $c_{\text{old}}$ is the cost of coding one traversal the long way
and $c_{\text{new}}$ the cost of naming it.
\end{proposition}

\begin{theorem}[The growth law]\label{thm:growthlaw}
\Proved\ Attach the concept if and only if $\Delta L_{\text{total}}<0$,
that is
\begin{equation}\label{eq:growthlaw}
  \boxed{\ n\;>\;
  \frac{b\,d_w\bigl(1+\sum_{e\in\gamma}d_e\bigr)}
       {c_{\text{old}}-c_{\text{new}}}\ }
\end{equation}
\end{theorem}
\begin{proof}
Add \eqref{eq:modelterm} and \eqref{eq:dataterm} and require the sum
negative; rearrange, using $c_{\text{old}}>c_{\text{new}}$.
\end{proof}

\begin{intu}\Intu
Read what this is: a \emph{recurrence threshold that is computed, not
chosen}. It is the ratio of what a concept costs to store against what
it saves per use --- the number of times a pattern must repeat before it
is worth naming.
\end{intu}

Three consequences, two of which were not asked for:

\begin{enumerate}
  \item $n$ is a number the engine can print, rather than a parameter a
        human sets.
  \item \textbf{It orders the holes.} When several harmonic classes
        compete, fill them in decreasing $\abs{\Delta L_{\text{total}}}$.
        This removes a decision that had been silently delegated to a
        human.
  \item \textbf{It predicts non-attachment.} A one-off inconsistency is
        \emph{noise to be tolerated}, not a concept to be invented. The
        earlier threshold-based law had no way to express this, which is
        exactly how it would have over-generated.
\end{enumerate}

\chapter{The Coned Concept}\label{ch:cone}

This is the central construction of the later work: the stalk of a
newly-grown concept, \emph{derived} rather than chosen.

\section{The topological cone}

\begin{construction}[Coning a cycle]\label{con:conetopo}
Let $\gamma=(v_1,e_1,v_2,\dots,v_k,e_k,v_1)$ be a $1$-cycle carrying
harmonic mass, with $e_i=\{v_i,v_{i+1}\}$ (indices mod $k$). Attach a
vertex $w$, edges $f_i=\{w,v_i\}$, and triangles
$t_i=\{w,v_i,v_{i+1}\}$.
\end{construction}

\begin{proposition}[The cone kills the cycle]\label{prop:conekills}
\Proved\ With $\partial t_i=e_i-f_{i+1}+f_i$,
\[
  \partial\Bigl(\sum_i t_i\Bigr)
  =\sum_i e_i+\sum_i(f_i-f_{i+1})=\gamma,
\]
since the $f$ terms telescope to zero. Hence $\gamma$ becomes a
boundary.
\end{proposition}

\section{The stalks, chosen minimally}

\begin{construction}[New edge stalks]\label{con:conestalks}
Set
\[
  \stalk{f_i}:=\stalk{v_i},
  \qquad
  \rest{v_i\triang f_i}:=\Id .
\]
\end{construction}

\begin{hon}\Hon
This is an \emph{engineering choice}, and a deliberately empty one. Any
distortion introduced here would be a free parameter, and whatever the
construction then achieved could be attributed to a cleverly chosen edge
stalk rather than to the concept. All the content is forced into
$\stalk{w}$ and the maps out of it, which is where it can be reasoned
about.
\end{hon}

The only remaining unknowns are $\stalk{w}$ and the legs
$p_i:=\rest{w\triang f_i}:\stalk{w}\to\stalk{v_i}$.

\section{The cone condition is forced, not imposed}

\begin{proposition}[Derivation of the cone condition]\label{prop:conecond}
\Proved\ Write $r_i^-:=\rest{v_i\triang e_i}$ and
$r_i^+:=\rest{v_{i+1}\triang e_i}$. For a $0$-cochain $x$:
\begin{enumerate}[label=(\roman*)]
  \item $\cob x$ vanishes on the new edges if and only if
        $x_{v_i}=p_i(x_w)$ for every $i$;
  \item given (i), $\cob x$ vanishes on the old edges for
        \emph{every} $x_w\in\stalk{w}$ if and only if
        \begin{equation}\label{eq:conecond}
          \boxed{\ r_i^+\circ p_{i+1}=r_i^-\circ p_i
          \quad\text{for all }i\ }
        \end{equation}
\end{enumerate}
\end{proposition}
\begin{proof}
(i) On the new edge $f_i$, using
Construction~\ref{con:conestalks},
\[
  (\cob x)_{f_i}
  =\rest{w\triang f_i}x_w-\rest{v_i\triang f_i}x_{v_i}
  =p_i(x_w)-x_{v_i},
\]
which vanishes exactly when $x_{v_i}=p_i(x_w)$.

(ii) Substituting $x_{v_i}=p_i(x_w)$ into the old edge $e_i$,
\[
  (\cob x)_{e_i}=r_i^+p_{i+1}(x_w)-r_i^-p_i(x_w)
  =\bigl(r_i^+p_{i+1}-r_i^-p_i\bigr)(x_w).
\]
A linear map vanishes on all of $\stalk{w}$ if and only if it is the
zero map, which is \eqref{eq:conecond}.
\end{proof}

\begin{intu}\Intu
Part (i) says something worth pausing on: \textbf{the value at $w$
generates the section over the whole cycle}. That is what it means for a
concept to explain a loop --- the loop stops being $k$ independent facts
and becomes one.
\end{intu}

Condition \eqref{eq:conecond} is verbatim the definition of a
\emph{cone} over the diagram
\[
  D_\gamma:\quad
  \stalk{v_1}\xrightarrow{r_1^-}\stalk{e_1}\xleftarrow{r_1^+}
  \stalk{v_2}\xrightarrow{r_2^-}\stalk{e_2}\xleftarrow{\ \ }\cdots
\]
with apex $\stalk{w}$ and legs $p_i$. The categorical cone condition was
not imposed on the construction; it fell out of demanding that the
topological cone do its job.

\section{The construction}

\begin{construction}[The coned concept]\label{con:coned}
\[
  \boxed{\ \stalk{w}:=\varprojlim D_\gamma,
  \qquad p_i:=\text{the limit's projections.}\ }
\]
\end{construction}

\begin{proposition}[Concrete description]\label{prop:conelim}
\Proved\ For finite-dimensional stalks,
\[
  \varprojlim D_\gamma
  =\Bigl\{(x_1,\dots,x_k)\in\bigoplus_i\stalk{v_i}\ \Big|\
    r_i^-x_i=r_i^+x_{i+1}\ \forall i\Bigr\}
  \;\cong\;H^0\bigl(\gamma;\Fs|_\gamma\bigr).
\]
\end{proposition}
\begin{proof}
The limit of a finite diagram of vector spaces is the subspace of the
product on which all the diagram's maps agree, which is the displayed
equaliser. Comparing with Definition~\ref{def:coboundary} restricted to
the subcomplex $\gamma$, the agreement conditions are exactly
$\cob x=0$, so the space is $\ker\cob|_\gamma=H^0(\gamma;\Fs|_\gamma)$.
\end{proof}

\begin{intu}\Intu
In words: the new concept \emph{is} everything that was consistent
around the loop but could not be glued. It does not summarise the
cycle's members; it \emph{is} the compatible part of them.
\end{intu}

Three properties, each of which had been written down separately as a
requirement:

\begin{center}
\begin{tabular}{@{}p{3.4cm}p{8.4cm}@{}}
\toprule
\textbf{Property} & \textbf{Why it matters}\\
\midrule
canonical --- unique up to unique isomorphism & no hand-made choice
enters the construction\\
path-dependent --- depends on $D_\gamma$, i.e. on which cycle history
produced & precisely the property whose absence caused an earlier
formal-concept-analysis approach to be rejected\\
degenerate case & if $H^0(\gamma;\Fs|_\gamma)=0$ the limit is the zero
space --- but see Proposition~\ref{prop:unreachable}\\
\bottomrule
\end{tabular}
\end{center}

\section{The refusal branch is unreachable}

\begin{proposition}[Harmonic mass implies sections]\label{prop:unreachable}
\Proved\ On a $k$-cycle with vertex and edge stalks all of dimension
$n$, and with no $2$-cells,
\[
  \dim H^1(\gamma;\Fs|_\gamma)=\dim H^0(\gamma;\Fs|_\gamma).
\]
Consequently the cycle carries harmonic mass if and only if it has
nonzero sections.
\end{proposition}
\begin{proof}
For a $1$-dimensional complex the Euler characteristic satisfies
$\chi=\dim\Czero-\dim\Cone=\dim H^0-\dim H^1$. On a $k$-cycle with all
stalks of dimension $n$ there are $k$ vertices and $k$ edges, so
$\dim\Czero=\dim\Cone=kn$ and $\chi=0$, giving the equality. With no
$2$-cells the harmonic space \emph{is} $H^1$.
\end{proof}

\begin{hon}\Hon
An earlier draft claimed the degenerate case is meaningful --- ``nothing
was consistent around that loop, so the construction refuses rather than
inventing''. True, and it never happens. By
Proposition~\ref{prop:unreachable}, $H^0=0$ implies no harmonic mass,
which implies no growth address, which implies the growth law never
fires on that cycle at all. Confirmed numerically: the degenerate case
reports that the harmonic space is already empty before coning, so there
is nothing to kill.

In general the gap is $\dim H^1-\dim H^0=k(n_e-n_v)$, so the branch
\emph{could} become reachable if edge stalks are ever given a different
dimension from vertex stalks. They are not, currently.
\end{hon}

\section{The limit is a ceiling, not a floor}\label{sec:conesize}

\begin{hon}\Hon
\Refuted\ An earlier draft justified
Construction~\ref{con:coned} by asserting that the limit is ``the
smallest object admitting the required maps''. That is backwards, and
the error is worth keeping visible.
\end{hon}

\begin{proposition}[Universality bounds from above]\label{prop:ceiling}
\Proved\ Every cone over $D_\gamma$ factors uniquely through the limit.
If a candidate apex $\stalk{w}'$ has jointly injective legs --- i.e.
carries no data that no restriction map ever sees --- then that
factoring map is injective, so
\[
  \dim\stalk{w}'\ \le\ \dim\varprojlim D_\gamma .
\]
\end{proposition}
\begin{proof}
The universal property of the limit gives a unique
$\phi:\stalk{w}'\to\varprojlim D_\gamma$ commuting with the legs. If
$\phi(z)=0$ then every leg $p_i'(z)=0$; joint injectivity of the legs
then forces $z=0$, so $\phi$ is injective and the dimension inequality
follows.
\end{proof}

So the limit is the \emph{largest non-redundant} choice, not the
smallest --- a ceiling on what a new concept can usefully carry. The two
determinations are therefore genuinely separate, and both are derived:

\begin{enumerate}
  \item \textbf{Consistency fixes the shape.} $\stalk{w}$ must be a cone
        over $D_\gamma$; the limit is the universal one.
  \item \textbf{MDL fixes the size.} By \eqref{eq:modelterm} the cost is
        linear in $d_w$, so the optimum is a \emph{subspace} of the
        limit: keep the directions that pay for themselves, discard the
        rest.
\end{enumerate}

\chapter{Retrieval: Rank, Do Not Threshold}\label{ch:retrieval}

\begin{intu}\Intu
Retrieval decides which stored concepts an organ sees. A threshold on
distance answers ``which concepts are close enough?'', which sounds
right and is not: it makes the number of results depend on how the
embedding space happens to be scaled, and it silently returns nothing at
all when the query lands in a sparse region.
\end{intu}

\begin{construction}[Bounded rank query]\label{con:topk}
Replace the $\varepsilon$-ball query with a bounded rank query returning
the top $k$ stored concepts, ranked by \emph{cosine similarity} to the
query mean, using a bounded max-heap of size $k$. Zero-norm vectors are
dropped rather than scored, since a zero vector has no direction.
\end{construction}

\begin{proposition}[Why cosine and not squared distance]\label{prop:cosine}
\Proved\ Expanding,
\[
  \norm{\mu_q-\mu_i}^2
  =\norm{\mu_q}^2-2\inner{\mu_q}{\mu_i}+\norm{\mu_i}^2 .
\]
For a fixed query, $\norm{\mu_q}^2$ is constant and drops out of any
ranking. The term $\norm{\mu_i}^2$ does \emph{not}. Hence squared
distance reproduces the cosine ranking if and only if every stored mean
is unit-norm.
\end{proposition}

\begin{hon}\Hon
Stored means are \emph{not} unit-norm. By
Construction~\ref{con:piv}, $\pi_v$ is a weighted centroid of unit
vectors, so its norm lies below $1$ and falls further the more spread
the concepts it summarises. Ranking by squared distance would therefore
apply a per-concept penalty \emph{proportional to how broad a concept
is}, pushing exactly the general concepts down the list for a reason
unrelated to the query.

There is also an evidential argument, which is decisive on its own: the
measured recall figure against which the retrieval fix is accepted was
obtained on a \emph{cosine} ranking. Ranking by squared distance is a
different ordering carrying an unmeasured bias, which would have made
the acceptance criterion meaningless.
\end{hon}

\begin{remark}[Cost]
Scoring needs only $\mu$; the low-rank factor $U$, the floor $D$ and the
stored reasoning text are not needed unless the row enters the heap.
Reading $\mu$, scoring, and deserialising the rest only on entry drops
the per-row cost from $O(dk^2+k^3)$ to $O(d)$ for the large majority of
rows.
\end{remark}

\begin{hon}\Hon
$k$ is an \emph{engineering choice} ($k=24$), not a derived quantity.
What changed is that it is now constrained on both sides --- too small
starves the organ, too large defeats the point of retrieval --- rather
than being an unexamined constant. Saying so is the difference between a
parameter and a magic number.
\end{hon}
